\section{Phần còn lại của danh mục}
\label{sec:ch15-phan-con-lai}

\index{mô hình nút thắt khái niệm!chín chỗ yếu|(}%
Câu mà mục trước khép lại là câu của chính người đề xuất khung, rằng
\tn{rò rỉ}{leakage} ở mô hình mềm thường không tránh được. Mục này đặt câu ấy
vào chỗ của nó, vì rò rỉ là
một trong chín chỗ yếu mà ngành tự liệt kê, và tám chỗ còn lại quyết định lối
thoát của Phần~V còn lại bao nhiêu.

Bản liệt kê là một khảo sát~\autocite{p31cbmrisks}. Nó không chạy thí nghiệm
nào, giống khảo sát biến thể LIME của
chương~\ref{ch:lime-nao-dang-tin}, nên mọi con số trong nó là con số của người
khác và giá trị của nó nằm ở cách sắp xếp. Nó chia theo hai nhánh mà
mục~\ref{sec:ch15-khai-niem-trong-mo-hinh} đã nêu: năm chỗ yếu ở nhánh có giám
sát và bốn ở nhánh không giám sát.

\index{mô hình nút thắt khái niệm!hiểu ngữ nghĩa hạn chế}%
Ở nhánh có giám sát, chỗ yếu thứ nhất là rò rỉ và bốn mục trên đã đọc nó. Chỗ
thứ hai là mô hình không thật sự hiểu nội dung ngữ nghĩa của khái niệm. Bằng
chứng gọn nhất cho điều đó là phép thử che: che đúng vùng ảnh chứa một khái
niệm, rồi xem mô hình có còn báo khái niệm ấy hiện diện hay không, và nó vẫn báo.
Cùng hướng ấy là quan sát rằng bản đồ nổi bật gắn với từng khái niệm không tập
trung vào vị trí của khái niệm mà lan cả sang nền. Chỗ yếu này khác rò rỉ ở đối
tượng: rò rỉ nói cột khái niệm mang thêm cái gì, còn ở đây cột khái niệm mang
đúng cái tên của nó nhưng đọc nó ra từ chỗ khác trong ảnh.

\index{mô hình nút thắt khái niệm!giả thiết khái niệm độc lập}%
Chỗ thứ ba là giả thiết các khái niệm độc lập với nhau. Giả thiết ấy tạm chấp
nhận được với khái niệm nhìn thấy trực tiếp như màu sắc, và hỏng với khái niệm
mang tính ngữ nghĩa như hình dạng, vì những khái niệm ấy chồng lấn nhau. Đây là
chỗ yếu duy nhất trong chín chỗ mà khảo sát ghi rằng chưa ai khảo sát có hệ
thống. Chỗ thứ tư là tính dễ vỡ trước phép nhiễu, và nó có hai mặt: nhiễu trong
không gian đầu vào, nơi đã có phép tấn công xóa một khái niệm, đưa vào một khái
niệm hoặc đánh tráo cả một tập khái niệm mà vẫn giữ nguyên nhãn dự đoán; và
nhiễu trong chính không gian khái niệm, nơi có phép cấy cửa hậu bằng cách thay
một tập khái niệm chuẩn bị sẵn để ép mô hình trả một nhãn định trước. Khảo sát
nhận xét rằng không gian khái niệm thưa hơn không gian đầu vào, dù nhỏ chiều
hơn, nên nó dễ bị tấn công hơn.

\index{mô hình nút thắt khái niệm!can thiệp không cân}%
Chỗ thứ năm đánh thẳng vào điều mà mục~\ref{sec:ch15-khai-niem-trong-mo-hinh} gọi
là toàn bộ sức hấp dẫn của họ mô hình này. Can thiệp được cho là chỗ con người
sửa được mô hình, và một nghiên cứu có hệ thống trên các mô hình lớn tìm ra rằng
có những tập khái niệm can thiệp vào thì không sửa được dự đoán sai, còn tập
khác thì sửa được dễ dàng, và điều đó xảy ra cả với những khái niệm có giá trị
rỗng. Kết luận của nghiên cứu ấy là can thiệp không phải một thao tác độc lập
đáng tin cậy mà là hệ quả trực tiếp của quy trình huấn luyện. Đọc cạnh
mục~\ref{sec:ch15-ro-ri} thì hai bên nói cùng một chuyện từ hai phía: ở đó can
thiệp thất bại vì nó xóa mất phần thông tin thêm mà đầu phân loại trông cậy, ở
đây nó thất bại vì quan hệ giữa các khái niệm học được không phải quan hệ mà
người can thiệp tưởng.

\index{mô hình nút thắt khái niệm!bốn chỗ yếu của nhánh không giám sát}%
Bốn chỗ yếu của nhánh không giám sát khác hẳn về hình dạng, vì ở đó không có
nhãn khái niệm nên không có mốc nào để lệch khỏi. Ba trong bốn nói được gọn, còn
chỗ thứ ba thì phải để riêng một đoạn vì cái tên nó mang. Chỗ thứ nhất là mô
hình học phải tương quan giả từ thiên lệch của dữ liệu, ví dụ trên một bộ ảnh
camera hành trình thì đặc trưng nền lấn át còn đèn đỏ và biển dừng bị bỏ qua.
Chỗ thứ hai là không có cách nào đọc thẳng một khái niệm đã học: người ta phải
suy ra nó từ những mẫu huấn luyện kích hoạt nó mạnh nhất, và việc hiểu tập mẫu
ấy nói lên điều gì thì để cho người dùng. Chỗ thứ tư là thế đánh đổi giữa hiệu
năng và tính diễn giải được đặc biệt gay ở nhánh này, vì thứ duy nhất ép khái
niệm phải đọc được là các phép chính quy hóa, và chính chúng làm hiệu năng tụt.

\index{tính ổn định!nghĩa thứ ba của chữ fidelity}%
Chỗ thứ ba mang một cái tên mà sách đã cấp cho thứ khác. Khảo sát gọi nó là
\enquote{concept fidelity} không nhất quán, và định nghĩa nó là việc các mẫu
cùng một lớp phải chia sẻ phần lớn khái niệm với nhau; điều quan sát được là các
phương pháp không giám sát học ra những khái niệm khác hẳn nhau giữa các lần
chạy, dù hiệu năng nhiệm vụ ngang nhau. Đây là lần thứ ba chữ ấy mang một nghĩa
khác trong sách, sau đại lượng fidelity của chương~\ref{ch:khung-hop-nhat} và
nhóm vấn đề mà chương~\ref{ch:lime-nao-dang-tin} gọi là độ khớp. Đọc theo
bảng~\ref{tab:ch01-bon-quan-he} thì nghĩa thứ ba này không phải hai nghĩa kia mà
là \tn{tính ổn định}{stability}: nó so hai lần chạy với nhau, không so một lời
giải thích với
đầu ra của mô hình và cũng không so nó với hành vi. Sách giữ chữ fidelity cho
đại lượng của chương~\ref{ch:khung-hop-nhat} và gọi chỗ yếu này bằng tên của
quan hệ mà nó thật sự nói tới.

\index{mô hình nút thắt khái niệm!không kiến trúc nào đủ sáu tính chất}%
Khảo sát khép phần có giám sát bằng một bảng đối chiếu, và bảng ấy là chỗ nó nói
được nhiều nhất. Mười kiến trúc, xếp theo thứ tự công bố từ CBM gốc
tới các biến thể mới nhất, đối chiếu với sáu tính chất: có chống rò rỉ không, có
mô hình hóa quá trình sinh dữ liệu theo lối xác suất không, có cấu trúc hóa
không gian khái niệm không, có ý thức về can thiệp không, có hiểu ngữ nghĩa
không, và có bền trước phép nhiễu không. Không kiến trúc nào đạt cả sáu, và
khảo sát nói thẳng điều đó trong chú thích bảng. Ba con số đọc được từ bảng.
Kiến trúc tốt nhất đạt bốn trong sáu. Cột cấu trúc hóa không gian khái niệm
chỉ có đúng một dấu tích trong mười dòng, khớp với nhận xét ở trên rằng đây là
chỗ yếu chưa ai khảo sát có hệ thống. Và CBM gốc, dòng đầu bảng,
không đạt tính chất nào trong sáu.

\index{mô hình nút thắt khái niệm!chín chỗ yếu|)}%
Còn hai điều về chính khảo sát mà chương phải nói. Thứ nhất, câu của nó về biểu
diễn cứng và biểu diễn mềm là câu sai chiều, và
mục~\ref{sec:ch15-dung-cu-duoc-kiem} đã đối chiếu với bài gốc rồi. Thứ hai, và
đáng chú ý hơn với sách này, mục hướng nghiên cứu tương lai của nó gồm ba hướng:
dùng mô hình ngôn ngữ và mô hình thị giác ngôn ngữ để phát hiện khái niệm, phòng
thủ trước các phép tấn công quyền riêng tư, và thiết kế kiến trúc tốt hơn. Cả ba
đều là hướng dựng mô hình. Một khảo sát dài mười trang về rủi ro và giới hạn của
cả một họ mô hình khái niệm, viết sau khi đã liệt kê chín chỗ yếu, không nêu một
hướng nào về việc đo. Bài của mục~\ref{sec:ch15-dung-cu-duoc-kiem} lên arXiv
trước khảo sát này và làm đúng việc ấy, nên chỗ trống này là chỗ trống trong tầm
nhìn của khảo sát chứ không phải trong ngành.
